\subsection{An obstruction to bonuses confined to saturated counts}\label{sec:unsaturated-obstruction}
Recovering the excess hidden by $b_i=6$ is one possible source of additional
charge. The completed-share graph has a minimum-mean cycle on which every
stored count is smaller than six. This prevents such bonuses alone from
improving its minimum mean. It is a direct instance of
Proposition~\ref{prop:bonus-criterion}: the old critical cycle remains in
the zero-bonus subgraph.

\begin{proposition}[An unchanged cycle under nonnegative bonuses]\label{prop:unsaturated-cycle}
Let a finite directed graph have edge weights $w$, minimum cycle mean 483,
and the following four-edge cycle. The count column records six consecutive
truncated counts on each edge.
\begin{center}
\begin{tabular}{c c r}
\toprule
Edge in cyclic order & Counts & Weight\\\midrule
1 & $(5,2,4,5,5,2)$ & 252\\
2 & $(2,4,5,5,2,4)$ & 480\\
3 & $(4,5,5,2,4,5)$ & 600\\
4 & $(5,5,2,4,5,5)$ & 600\\
\bottomrule
\end{tabular}
\end{center}
Replace $w(e)$ by $w(e)+B(e)$, where $B(e)\ge0$ on every edge and $B(e)=0$
whenever all six stored counts on that edge are smaller than six. Keep the
same graph. Its minimum cycle mean remains 483. No potential for these new
weights can satisfy~\eqref{eq:strict234} with positive scale and nonnegative
range.
\end{proposition}
\begin{proof}
Every cycle has nondecreasing total weight, so its mean remains at least 483.
The displayed cycle has unchanged total
\[
252+480+600+600=1932=4\cdot483.
\]
It still attains 483, proving equality of the minimum cycle means.

Suppose a potential satisfies $p+\pi(v)\le d(w(e)+B(e))+\pi(u)$ on every
edge $e=(u,v)$, with $d>0$ and range at most $R\ge0$. Summing over the displayed
cycle cancels its four potential differences and gives $p\le483d$. Therefore
\[
227p-R\le227\cdot483d=109641d<111960d,
\]
contradicting the strict criterion. The deficit at scale one and range zero
is already $111960-227\cdot483=2319$. The cycle arithmetic and monotonicity
of potential inequalities under nonnegative bonuses have passed Lean checks.
The cycle's membership in the saved graph and the original minimum mean were
verified separately against all graph edges~\cite{unsaturatedarchive}.
\end{proof}

The hypothesis concerns the current graph and the current global budget.
It does not assert that this cycle is realizable by an actual snake. Proving
additional geometric exclusions, restricting to realizable walks, or changing
the boundary estimate can invalidate its use as an obstruction. Bonuses on
edges whose counts are all below six also lie outside the stated restriction.

The same witness identifies two relevant local targets. The proposed floor
20 for $012034$ applies twice, and the proposed floor 16 for $012340$ applies
once. With both floors, its total would rise by 56 to 1988, giving mean 497.
This calculation concerns only the displayed cycle; another cycle may have
a lower mean in the modified graph. Both local targets remain unproved.
Together with the single-target comparison above, this explains why completing
the $012034$ proof alone cannot settle the stronger bound.
